% appF_tail.tex: proof of Proposition prop:tail (localisation of the spectral weights).

\section{Proof of the localisation bound}\label{app:tail}

This appendix proves Proposition~\ref{prop:tail}. The spectral weights $m_n(k)$ form a probability measure on the
rescaled nodes, and the mass it gives to a short interval $[0,r]$ at the hard edge is controlled by its Stieltjes
transform at one point off the contour (Lemma~\ref{lem:stieltjes}). That transform is the product of two entries of the first row of
$Q$ (Lemma~\ref{lem:rowproduct}), so the bound reduces to an estimate of $Q$ at a single point at distance $\asymp r$
from all contours: in case~(a) from the comparison argument of Section~\ref{sec:hardedge}, in case~(b) from the error
matrix of \S\ref{sec:classical-small} and explicit bounds on Bessel functions. The notation is that of
\S\ref{sec:tail}.

\begin{lemma}[Stieltjes bound]\label{lem:stieltjes}
Let $\nu$ be a probability measure on $[0,\infty)$, $r>0$ and $z=re^{5\pi i/6}$. Then
$\nu([0,r])\le\frac8{\sqrt3}\,r\,\bigl|\int\diff\nu(X)/(X-z)\bigr|$.
\end{lemma}

\begin{proof}
For $X\ge0$, $\re\frac1{X-z}=\frac{X+r|\cos(5\pi/6)|}{|X-z|^2}\ge0$, and for $0\le X\le r$ also $|X-z|^2\le4r^2$, so
$\re\frac1{X-z}\ge\frac{\sqrt3/2}{4r}$. Hence $\nu([0,r])\frac{\sqrt3}{8r}\le\re\int\frac{\diff\nu(X)}{X-z}$.
\end{proof}

\begin{lemma}[Row product]\label{lem:rowproduct}
Let $\nu_n:=\sum_km_n(k)\delta_{X_k}$, with the nodes $X_k=(n/k)^p$ of~\eqref{eq:Xk}, and $S_n(z):=\int\diff\nu_n(X)/(X-z)$.
Suppose that $Q=\mathsf c\,U\,\mathsf k^{\sigma_3}$ on a connected open set containing $z\notin\{X_k\}\cup\{0\}$ and a
neighbourhood of infinity, where $\mathsf c=\mathrm{diag}(\mathsf c_1,\mathsf c_2)$ is constant, $\mathsf k$ is scalar, and
$Q=I+Q_1/z+O(z^{-2})$. Then $S_n(z)=-Q_{11}(z)Q_{12}(z)/(Q_1)_{12}$.
\end{lemma}

\begin{proof}
By~\eqref{eq:U-sum}, $U_{12}(z)=\sum_kn^pw_kU_{11}(X_k)/(z-X_k)$, and $U_{11}(X)=n^{pn}P_n(Xn^{-p})$ is monic of degree $n$
in $X$. Since $(U_{11}(z)-U_{11}(X))/(z-X)$ is a polynomial of degree $n-1$ in $X$, orthogonality gives
$U_{11}(z)U_{12}(z)=\sum_kn^pw_kU_{11}(X_k)^2/(z-X_k)=-\mathsf hS_n(z)$ with $\mathsf h:=\sum_kn^pw_kU_{11}(X_k)^2$, because
$m_n(k)=n^pw_kU_{11}(X_k)^2/\mathsf h$. Expanding $U_{12}$ at infinity, orthogonality kills the terms $z^{-j-1}$, $j<n$, and
the coefficient of $z^{-n-1}$ is $\mathsf h$; so the $(1,2)$ entry of the first coefficient of $Uz^{-n\sigma_3}$ is $\mathsf h$.
Now $Q_{11}Q_{12}=\mathsf c_1^2U_{11}U_{12}$, and $Q_{11}\to1$ forces $\mathsf c_1z^n\mathsf k(z)\to1$, so
$(Q_1)_{12}=\lim_{z\to\infty}zQ_{12}=\mathsf c_1^2\mathsf h$.
\end{proof}

The identity of Lemma~\ref{lem:rowproduct} expresses the diagonal of the resolvent of the Jacobi matrix through the
product of the two solutions of the recurrence; we use it only in the form stated.

\begin{proof}[Proof of Proposition~\ref{prop:tail}]
\emph{Reduction.} Take $j=n$ first. Since $k>Mn$ exactly when $X_k<M^{-p}$, we have
$\sum_{k>Mn}m_n(k)\le\nu_n([0,r])$ with $r:=M^{-p}\in(0,1]$. Let $z:=re^{5\pi i/6}$. For $r\le\frac14\delta$ the point
$z$ lies in $\Oinf$ (the lips stay in $|\arg z|\le\frac54\thlip<\pi/2$) and off all contours, and there
$Q=\mathsf c\,U(e^{n\varphi}FD_W)^{\sigma_3}$ by~\eqref{eq:VM} and~\eqref{eq:Q}, as also near infinity; so
Lemmas~\ref{lem:stieltjes} and~\ref{lem:rowproduct} give
\begin{equation}\label{eq:tail-reduction}
  \nu_n([0,r])\le\frac8{\sqrt3}\,r\,\frac{|Q_{11}(z)Q_{12}(z)|}{|(Q_1)_{12}|},\qquad (Q_1)_{12}=ic+O(1/n),
\end{equation}
the latter by Proposition~\ref{prop:moments} in case~(a) and by Proposition~\ref{prop:classical} in case~(b), together
with~\eqref{eq:P1P2}. For $r_0<r\le1$, where $r_0:=\frac14\delta$ in case~(a) and $r_0:=\frac12\delta_1$ in case~(b)
($\delta_1$ is fixed below), we use $\nu_n([0,r])\le1\le r_0^{-1/2}r^{1/2}$. For $j=n-1$ apply the
result at degree $n-1$ with $M':=Mn/(n-1)\ge M$; this costs a factor $\frac12$ in $\gamma_*$. The geometric fact used below is
that for $s'$ on the band or on a lip in $\mathbb D(0,\delta)$ the angle between $s'$ and $z$ is at least $\pi/3$, so
$|s'-z|\ge\frac12(|s'|+|z|)$; and that $|f(z)|\asymp r$, $\zeta_z:=n^2f(z)$ satisfies $|\arg\zeta_z|<\pi/6+\pi/24$.

\emph{Case (a).} Here $1<s<1+2p$, $|\alpha|<1$, $\ell\equiv1$, so $\varepsilon_\Lambda\equiv0$ and
$\varepsilon'_n=n^{-2\nu_p}\log n$. Put $y:=n^2r$, $D_z:=(2\pi n)^{\sigma_3/2}\zeta_z^{-\sigma_3/4}$ and let $y_0\ge1$ be such
that $|\zeta_z|\ge2$ for $y\ge y_0$. We show
\begin{equation}\label{eq:Q-ray}
  |Q(z)|\le Cr^{-1/4}\qquad(y_0n^{-2}\le r\le\tfrac14\delta),
\end{equation}
which with~\eqref{eq:tail-reduction} gives $\nu_n([0,r])\le Cr^{1/2}=CM^{-p/2}$ for $1\le M\le(n^2/y_0)^{1/p}$.

By Proposition~\ref{prop:UR-comparison}(i), Proposition~\ref{prop:UR-true} and \S\ref{sec:tool-cd}(b), as in the proof of
Proposition~\ref{prop:moments}, $\mathsf K:=Q\tilde Q^{-1}=I+\Cauchy(\mathcal E)$ off $\hat\Sigma$, with
$\mathcal E=Q_-(v-\tilde v)\tilde v^{-1}\tilde Q_-^{-1}$. For $|\zeta_z|\ge2$, $\tilde Q(z)=A_0(z)D_z\Theta_\infty(\zeta_z)$ by
\eqref{eq:he-P0hat} and~\eqref{eq:he-Phi-large}, and $|A_0(z)D_z|\le Cr^{-1/4}$ since the entries of $D_z$ have moduli
$\asymp r^{\mp1/4}$. Since $Q=[A_0D_z+(\mathsf K-I)A_0D_z]\Theta_\infty$, it suffices to show
\begin{equation}\label{eq:K-weighted}
  |(\mathsf K(z)-I)A_0(z)D_z|\le C(\varepsilon'_n+n^{-1})\,r^{-1/4}.
\end{equation}
The right factor $A_0(z)D_z$ is the scale of the comparison solution at $z$; without it the bound would not be small at
the Bessel scale. We split $(v-\tilde v)\tilde v^{-1}=\Delta_1+\Delta_2+\Delta_3$ and $Q_-=\tilde Q_-+(Q_--\tilde Q_-)$ as in
the proof of Proposition~\ref{prop:moments}, and put $\mathfrak m_n(s'):=|s'|^{\nu_p}\log(2/|s'|)$, so that
$|d_n(s')|\le C\mathfrak m_n(s')e^{-n|s'|^{1/2}/C}$ by~\eqref{eq:he-mismatch} and~\eqref{eq:he-localisation}. Write
$\varrho:=|s'|$, $t:=n\varrho^{1/2}$ and $B(s'):=A_0(s')^{-1}A_0(z)$; since $A_0^{\pm1}$ are bounded and $A_0$ is Lipschitz on
$\mathbb D(0,\delta/2)$, $|B(s')-I|\le C(\varrho+r)$.

\emph{The lips, $|\zeta_{s'}|\ge2$.} The integrand is the rank-one matrix
$d_n(\tilde Q(s')e_2)(e_1^{\mathsf T}\Theta_\infty(\zeta_{s'})^{-1}N)$, $N:=D_{s'}^{-1}BD_z$. Here $|\tilde Q(s')e_2|\le
C\varrho^{-1/4}$, and the entries of $D_w^{\pm1}$ have moduli $\asymp|w|^{\mp1/4},|w|^{\pm1/4}$, which with $|B-I|\le C(\varrho+r)$
gives $|N|\le C((\varrho/r)^{1/4}+(r/\varrho)^{1/4})$. So, using $|\diff s'|\le C\diff\varrho$ and $|s'-z|\ge\frac12(\varrho+r)$, this
part of~\eqref{eq:K-weighted} is at most
\begin{align*}
  &Cr^{-1/4}\int\mathfrak m_ne^{-t/C}\bigl[1+(r/\varrho)^{1/2}\bigr]\frac{\diff\varrho}{\varrho+r}\\
  &\qquad=Cr^{-1/4}\int_{t\ge1/C}\mathfrak m_n\bigl(\tfrac{t^2}{n^2}\bigr)e^{-t/C}\Bigl[\frac{2t}{t^2+y}+\frac{2\sqrt y}{t^2+y}\Bigr]\diff t
  \le C\varepsilon'_nr^{-1/4},
\end{align*}
since $2t\le t^2+y$ (as $y\ge1$), $2\sqrt y\le(t^2+y)/t$, and $\mathfrak m_n(t^2/n^2)\le Ct^{2\nu_p}n^{-2\nu_p}(\log n+|\log t|+1)$.

\emph{The lips, $|\zeta_{s'}|\le2$.} Then $\varrho\le Cn^{-2}\le Cr$ and $\mathfrak m_n(s')\le C\varepsilon'_n$. By
\eqref{eq:he-Phi-small} and Lemma~\ref{lem:bessel-local}(iii), $|\tilde Q(s')e_2|\le Cn^{1/2}|\zeta_{s'}|^{-|\alpha|/2}$, and
\[
  e_1^{\mathsf T}\tilde Q(s')^{-1}A_0(z)D_z=e_1^{\mathsf T}\Phi_\alpha(\zeta_{s'})^{-1}(2\pi n)^{-\sigma_3/2}B(2\pi n)^{\sigma_3/2}
  \zeta_z^{-\sigma_3/4},
\]
whose norm is at most $C|\zeta_{s'}|^{-|\alpha|/2}|\zeta_z|^{1/4}\le C|\zeta_{s'}|^{-|\alpha|/2}n^{1/2}r^{1/4}$
(the lower-left entry of $B$ is $O(r)$ and $nr|\zeta_z|^{-1/4}\asymp|\zeta_z|^{1/4}r^{1/2}$). With $|s'-z|\ge r/2$ and
$\int_0^{Cn^{-2}}(n^2\varrho)^{-|\alpha|}\diff\varrho\le Cn^{-2}$ ($|\alpha|<1$; logarithms at $\alpha=0$), this part is at most
$C\varepsilon'_nn\,r^{1/4}r^{-1}n^{-2}=C\varepsilon'_nr^{-1/4}y^{-1/2}$.

\emph{The circles and the rest.} On $\Sigma_\circ$, $|\Delta_2|\le C/n$, $|\tilde Q_-^{\pm1}|\le C$ and $|s'-z|\ge\frac34\delta$,
which gives $Cn^{-1}r^{-1/4}$. On the lips, $|E|\le Ce^{-n|s'|^{-1/p}/C}$ and $|\tilde Q_-^{\pm1}|\le Cn^{1/2}|s'|^{-1/2}$, which
gives $Cr^{-1/4}e^{-n/C}$ (as $r\ge y_0n^{-2}$); the rest of $\hat\Sigma$, of length $O(n^p)$, gives $Cn^pe^{-n/C}r^{-1/4}$.

\emph{The correction term.} By Cauchy--Schwarz and~\eqref{eq:he-L2-diff} it is at most
$C(\varepsilon'_nn^{-1/2}+n^{-1})\,\|(v-\tilde v)\tilde v^{-1}\tilde Q_-^{-1}A_0(z)D_z/(s'-z)\|_{L^2(\hat\Sigma)}$. In the last norm, the
part $|\zeta_{s'}|\le2$ of $\Delta_1$ is at most $C\varepsilon'_nn^{1/2}r^{1/4}r^{-1}n^{-1}=C\varepsilon'_nn^{-1/2}r^{-3/4}$; the part
$|\zeta_{s'}|\ge2$ has square at most
\begin{align*}
  &C\int\mathfrak m_n^2e^{-2t/C}\Bigl[\Bigl(\frac\varrho r\Bigr)^{1/2}+\Bigl(\frac r\varrho\Bigr)^{1/2}\Bigr]\frac{\diff\varrho}{(\varrho+r)^2}\\
  &\qquad=C\int\mathfrak m_n^2e^{-2t/C}\Bigl(\frac{2t^2}{\sqrt y}+2\sqrt y\Bigr)\frac{n^2\diff t}{(t^2+y)^2}\le C\varepsilon_n'^2n^2y^{-3/2},
\end{align*}
that is, the part is at most $C\varepsilon'_nn^{-1/2}r^{-3/4}$; $\Delta_2$ gives $Cn^{-1}r^{-1/4}$ and $\Delta_3$ an exponentially small
multiple of $r^{-1/4}$. Since $nr^{1/2}=\sqrt y\ge1$, the product is at most $Cr^{-1/4}(\varepsilon'_n+n^{-1})$. Summing the
four parts gives~\eqref{eq:K-weighted}, hence~\eqref{eq:Q-ray}, and~\eqref{eq:tail-arcsine} follows with $\gamma_*=\frac12y_0^{-1/p}$.

\emph{Case (b).} Here $p>1$ and $\alpha>-1$ is arbitrary, and we work in the frame of \S\ref{sec:classical-small}.
First, $|\mathsf R(z)-I|\le C/n$ for $0<r\le\frac12\delta_1$, with a fixed $\delta_1\le\delta$. To see this, note
$\mathsf R(z)-I=\Cauchy(\mathsf R_-(J_{\mathsf R}-I))(z)$ by Proposition~\ref{prop:classical}(iii); the part of
$\Sigma_{\mathsf R}$ outside $\mathbb D(0,\delta_1)$ is at distance $\ge\frac12\delta_1$ from $z$ and contributes $O(1/n)$
by Proposition~\ref{prop:classical}(ii), (iv); on $\Sigma_{\mathsf R}\cap \mathbb D(0,\delta_1)$ (two lips leaving $0$ at angles
$\pm\pi/3$ and the segment $(-\delta_1,0)$) we have $|s'-z|\ge\frac14(|s'|+|z|)$ for $\delta_1$ small, and by
Lemma~\ref{lem:mismatch-classical} and Cauchy--Schwarz the contribution is $O(n^{-K})$. Next, in $U_0$,
$Q=\mathsf R\check A(2\pi n)^{\sigma_3/2}\sigma_3\Psi_\alpha(\zeta)\sigma_3\,(e^{-2\zeta^{1/2}}\rho_n^{-1})^{\sigma_3}$, and the last
factor is diagonal with determinant $1$, so it does not change $Q_{11}Q_{12}$. With
$(\mathsf r_1,\mathsf r_2):=e_1^{\mathsf T}\mathsf R(z)\check A(z)$, which is bounded by Lemma~\ref{lem:classical-param}(i),
and $u:=2\zeta_z^{1/2}$, \cite[(6.23)]{KMVV} gives
\[
  Q_{11}=\sqrt{2\pi n}\Bigl[\mathsf r_1I_\alpha(u)-\frac{i\,\mathsf r_2}{2n}uI_\alpha'(u)\Bigr],\qquad
  Q_{12}=-\sqrt{2\pi n}\Bigl[\frac{i\,\mathsf r_1}\pi K_\alpha(u)+\frac{\mathsf r_2}{2\pi n}uK_\alpha'(u)\Bigr].
\]
Here $|\mathrm{ph}\,u|\le\pi/3$. We use the following bounds, for $|\mathrm{ph}\,u|\le\pi/3$ and with $C$ depending on $\alpha$:
$|I_\alpha(u)|+|uI_\alpha'(u)|\le C|u|^\alpha$, $|K_\alpha(u)|\le C(1+\log(4/|u|))|u|^{-|\alpha|}$ and $|uK_\alpha'(u)|\le
C|u|^{-|\alpha|}$ for $|u|\le2$; and, with $\mathfrak r:=\re u\ge|u|/2$, $|I_\alpha(u)|+|I_\alpha'(u)|\le C(\mathfrak r^{-1/2}e^{\mathfrak r}+1)$
and $|K_\alpha(u)|+|K_\alpha'(u)|\le C\mathfrak r^{-1/2}e^{-\mathfrak r}$ for $\mathfrak r\ge1$. The first pair follows from the series
(10.25.2), with $\Gamma(\alpha+k+1)$ bounded below for $\alpha>-1$; the others from (10.32.9),
$K_\alpha(u)=\int_0^\infty e^{-u\cosh t}\cosh(\alpha t)\diff t$, which gives $|K_\alpha(u)|\le K_\alpha(\mathfrak r)$ and the same for
the derivative, and from (10.32.4) for $I_\alpha$, by elementary estimates of the integrals (for small $\mathfrak r$ use
$\cosh t\ge e^t/2$, for large $\mathfrak r$ use $\cosh t\ge1+t^2/2$ and $1-\cos\theta\ge2\theta^2/\pi^2$). Since $|u|\le Cn$, these
give $|Q_{11}Q_{12}|\le Cn\,b(|u|)$ with $b(\varrho):=(1+\log(4/\varrho))\varrho^{\alpha-|\alpha|}$ for $\varrho\le2$ and $b(\varrho):=1/\varrho$
for $\varrho\ge2$. By~\eqref{eq:tail-reduction}, with $|u|\asymp\sqrt y$, $y=n^2r$:
if $|u|\ge2$ then $\nu_n([0,r])\le Crn/(n\sqrt r)=Cr^{1/2}$; if $|u|\le2$ then
$\nu_n([0,r])\le Cn^{-1}y^{1+\min(\alpha,0)}(1+\log(C/y))$. With $r=M^{-p}$ and $p(1+\min(\alpha,0))=\min(p,s-1)=:\beta$, the
first case gives~\eqref{eq:tail-arcsine} for $M\le Cn^{2/p}$; the second gives it for the remaining
$M\le \gamma_*n^{2/p}$, where $y$ is bounded below and above, so that $n^{-1}\le CM^{-p/2}$; and the second gives
$\sum_{k>Mn}m_n(k)\le Cn^{-1}(M/n^{2/p})^{-\beta}(1+\log(M/n^{2/p}))$ for $M\ge n^{2/p}$. For $\tau<\min(p/2,s-1)\le\beta$ we have
$n^{-1}=(n^{2/p})^{-p/2}\le(n^{2/p})^{-\tau}$ and $(M/n^{2/p})^{-\beta}(1+\log(M/n^{2/p}))\le C_\tau(M/n^{2/p})^{-\tau}$, so the
second bound is at most $C_\tau M^{-\tau}$; and $M^{-p/2}\le M^{-\tau}$ in the first range.
\end{proof}
